% ============================================================================
% PHỤ LỤC: MÃ NGUỒN ĐẦY ĐỦ
% ============================================================================
\appendix
\chapter{Mã nguồn đầy đủ}
\label{chap:appendix}

\section{CUDA Kernels}

\subsection{Vector Sum (vecto.cu)}

\begin{lstlisting}[caption={Vector Sum Kernel hoàn chỉnh}]
#include <cuda_runtime.h>
#include <stdio.h>

#define N 1000000
#define BLOCK_SIZE 256

__global__ void vectorSum(float *a, float *b, float *c, int n) {
    int idx = blockIdx.x * blockDim.x + threadIdx.x;
    if (idx < n) {
        c[idx] = a[idx] + b[idx];
    }
}

int main() {
    float *h_a, *h_b, *h_c;
    float *d_a, *d_b, *d_c;
    
    // Allocate host memory
    h_a = (float*)malloc(N * sizeof(float));
    h_b = (float*)malloc(N * sizeof(float));
    h_c = (float*)malloc(N * sizeof(float));
    
    // Initialize
    for (int i = 0; i < N; i++) {
        h_a[i] = 1.0f;
        h_b[i] = 2.0f;
    }
    
    // Allocate device memory
    cudaMalloc(&d_a, N * sizeof(float));
    cudaMalloc(&d_b, N * sizeof(float));
    cudaMalloc(&d_c, N * sizeof(float));
    
    // Copy to device
    cudaMemcpy(d_a, h_a, N * sizeof(float), cudaMemcpyHostToDevice);
    cudaMemcpy(d_b, h_b, N * sizeof(float), cudaMemcpyHostToDevice);
    
    // Launch kernel
    int gridSize = (N + BLOCK_SIZE - 1) / BLOCK_SIZE;
    vectorSum<<<gridSize, BLOCK_SIZE>>>(d_a, d_b, d_c, N);
    
    // Copy result back
    cudaMemcpy(h_c, d_c, N * sizeof(float), cudaMemcpyDeviceToHost);
    
    // Verify
    printf("c[0] = %f, c[N-1] = %f\n", h_c[0], h_c[N-1]);
    
    // Cleanup
    free(h_a); free(h_b); free(h_c);
    cudaFree(d_a); cudaFree(d_b); cudaFree(d_c);
    
    return 0;
}
\end{lstlisting}

\subsection{Parallel Reduction (parallel\_reduction.cu)}

\begin{lstlisting}[caption={Parallel Reduction - 3 Methods}]
#include <cuda_runtime.h>
#include <stdio.h>

#define BLOCK_SIZE 256

// Method 1: Neighboring pairs (divergent warps)
__global__ void reduceNeighbors(int *input, int *output, int n) {
    __shared__ int sdata[BLOCK_SIZE];
    
    int tid = threadIdx.x;
    int idx = blockIdx.x * blockDim.x + threadIdx.x;
    
    sdata[tid] = (idx < n) ? input[idx] : 0;
    __syncthreads();
    
    for (int stride = 1; stride < blockDim.x; stride *= 2) {
        if (tid % (2 * stride) == 0) {
            sdata[tid] += sdata[tid + stride];
        }
        __syncthreads();
    }
    
    if (tid == 0) output[blockIdx.x] = sdata[0];
}

// Method 2: Improved - sequential addressing
__global__ void reduceNeighbors2(int *input, int *output, int n) {
    __shared__ int sdata[BLOCK_SIZE];
    
    int tid = threadIdx.x;
    int idx = blockIdx.x * blockDim.x + threadIdx.x;
    
    sdata[tid] = (idx < n) ? input[idx] : 0;
    __syncthreads();
    
    for (int stride = 1; stride < blockDim.x; stride *= 2) {
        int index = 2 * stride * tid;
        if (index < blockDim.x) {
            sdata[index] += sdata[index + stride];
        }
        __syncthreads();
    }
    
    if (tid == 0) output[blockIdx.x] = sdata[0];
}

// Method 3: Interleaved - no bank conflicts
__global__ void reduceInterleaved(int *input, int *output, int n) {
    __shared__ int sdata[BLOCK_SIZE];
    
    int tid = threadIdx.x;
    int idx = blockIdx.x * blockDim.x + threadIdx.x;
    
    sdata[tid] = (idx < n) ? input[idx] : 0;
    __syncthreads();
    
    for (int stride = blockDim.x / 2; stride > 0; stride >>= 1) {
        if (tid < stride) {
            sdata[tid] += sdata[tid + stride];
        }
        __syncthreads();
    }
    
    if (tid == 0) output[blockIdx.x] = sdata[0];
}
\end{lstlisting}

\subsection{Matrix Transpose}

\begin{lstlisting}[caption={Matrix Transpose với Shared Memory}]
#define TILE_DIM 32

// Naive implementation
__global__ void transposeNaive(float *in, float *out, int rows, int cols) {
    int ix = blockIdx.x * blockDim.x + threadIdx.x;
    int iy = blockIdx.y * blockDim.y + threadIdx.y;
    
    if (ix < cols && iy < rows) {
        out[ix * rows + iy] = in[iy * cols + ix];
    }
}

// Optimized with shared memory
__global__ void transposeShared(float *in, float *out, int rows, int cols) {
    __shared__ float tile[TILE_DIM][TILE_DIM + 1];
    
    int x = blockIdx.x * TILE_DIM + threadIdx.x;
    int y = blockIdx.y * TILE_DIM + threadIdx.y;
    
    if (x < cols && y < rows) {
        tile[threadIdx.y][threadIdx.x] = in[y * cols + x];
    }
    __syncthreads();
    
    x = blockIdx.y * TILE_DIM + threadIdx.x;
    y = blockIdx.x * TILE_DIM + threadIdx.y;
    
    if (x < rows && y < cols) {
        out[y * rows + x] = tile[threadIdx.x][threadIdx.y];
    }
}
\end{lstlisting}

\section{C\# Parallel Code}

\subsection{Factorial với Parallel.For}

\begin{lstlisting}[style=csharpstyle, caption={FactorialParallel.cs}]
using System;
using System.Threading.Tasks;

namespace Session5 {
    public class FactorialParallel {
        public static void Run() {
            int n = 10;
            
            Parallel.For(0, n, i => {
                long factorial = ComputeFactorial(i);
                Console.WriteLine($"{i}! = {factorial}");
            });
        }
        
        public static long ComputeFactorial(int num) {
            long result = 1;
            for (int j = 1; j <= num; j++) {
                result *= j;
            }
            return result;
        }
    }
}
\end{lstlisting}

\subsection{Vector Sum với Local State}

\begin{lstlisting}[style=csharpstyle, caption={VectorSumParallel.cs}]
using System;
using System.Linq;
using System.Threading;
using System.Threading.Tasks;

namespace Session5 {
    public class VectorSumParallel {
        public static void Run() {
            int vectorSize = 1000000;
            int[] vector = Enumerable.Range(1, vectorSize).ToArray();
            
            // Parallel with Local State
            long sum = 0;
            object lockObj = new object();
            
            Parallel.For(0, vectorSize,
                () => 0L,
                (i, state, localSum) => localSum + vector[i],
                (localSum) => {
                    lock (lockObj) {
                        sum += localSum;
                    }
                }
            );
            
            Console.WriteLine($"Sum = {sum}");
        }
    }
}
\end{lstlisting}

\subsection{Channel Pattern}

\begin{lstlisting}[style=csharpstyle, caption={FactorialChannel.cs}]
using System;
using System.Threading.Channels;
using System.Threading.Tasks;

namespace Session5 {
    public class FactorialChannel {
        public static async Task RunAsync() {
            int n = 10;
            var channel = Channel.CreateUnbounded<(int, long)>();
            
            // Producer
            var producer = Task.Run(async () => {
                for (int i = 0; i < n; i++) {
                    long fact = ComputeFactorial(i);
                    await channel.Writer.WriteAsync((i, fact));
                }
                channel.Writer.Complete();
            });
            
            // Consumer
            await foreach (var (num, fact) in channel.Reader.ReadAllAsync()) {
                Console.WriteLine($"{num}! = {fact}");
            }
            
            await producer;
        }
        
        public static long ComputeFactorial(int num) {
            long result = 1;
            for (int j = 1; j <= num; j++) result *= j;
            return result;
        }
    }
}
\end{lstlisting}

\section{Dot Product C\#}

\begin{lstlisting}[style=csharpstyle, caption={DotProduct với Task và Channel}]
using System;
using System.Linq;
using System.Threading.Channels;
using System.Threading.Tasks;

public class DotProductCS {
    // Sequential
    public static double Sequential(double[] A, double[] B) {
        double sum = 0.0;
        for (int i = 0; i < A.Length; i++)
            sum += A[i] * B[i];
        return sum;
    }
    
    // Task-based
    public static double WithTasks(double[] A, double[] B) {
        int numChunks = Environment.ProcessorCount;
        int chunkSize = (A.Length + numChunks - 1) / numChunks;
        double[] partialSums = new double[numChunks];
        Task[] tasks = new Task[numChunks];
        
        for (int c = 0; c < numChunks; c++) {
            int chunk = c;
            int start = chunk * chunkSize;
            int end = Math.Min(start + chunkSize, A.Length);
            
            tasks[c] = Task.Run(() => {
                double local = 0.0;
                for (int i = start; i < end; i++)
                    local += A[i] * B[i];
                partialSums[chunk] = local;
            });
        }
        
        Task.WaitAll(tasks);
        return partialSums.Sum();
    }
    
    // Channel-based
    public static async Task<double> WithChannel(double[] A, double[] B) {
        int numChunks = Environment.ProcessorCount;
        int chunkSize = (A.Length + numChunks - 1) / numChunks;
        var channel = Channel.CreateUnbounded<double>();
        
        var producers = new Task[numChunks];
        for (int c = 0; c < numChunks; c++) {
            int start = c * chunkSize;
            int end = Math.Min(start + chunkSize, A.Length);
            
            producers[c] = Task.Run(async () => {
                double local = 0.0;
                for (int i = start; i < end; i++)
                    local += A[i] * B[i];
                await channel.Writer.WriteAsync(local);
            });
        }
        
        _ = Task.Run(async () => {
            await Task.WhenAll(producers);
            channel.Writer.Complete();
        });
        
        double total = 0.0;
        await foreach (var partial in channel.Reader.ReadAllAsync()) {
            total += partial;
        }
        return total;
    }
}
\end{lstlisting}

\section{Cách biên dịch và chạy}

\subsection{CUDA}

\begin{lstlisting}[language=bash, caption={Biên dịch và chạy CUDA}]
# Bien dich
nvcc -o vector_sum vector_sum.cu

# Chay
./vector_sum

# Voi toi uu
nvcc -O3 -arch=sm_75 -o optimized program.cu

# Debug
nvcc -G -g -o debug program.cu
cuda-gdb ./debug
\end{lstlisting}

\subsection{C\#}

\begin{lstlisting}[language=bash, caption={Biên dịch và chạy C\#}]
# Tao project
dotnet new console -n MyProject

# Build
dotnet build

# Chay
dotnet run

# Release build
dotnet build -c Release
\end{lstlisting}

\section{Cấu hình Google Colab cho CUDA}

\begin{lstlisting}[language=python, caption={Setup CUDA trên Google Colab}]
# Kiem tra GPU
!nvidia-smi

# Viet file .cu
%%writefile program.cu
#include <stdio.h>
__global__ void hello() {
    printf("Hello from thread %d\n", threadIdx.x);
}
int main() {
    hello<<<1, 10>>>();
    cudaDeviceSynchronize();
    return 0;
}

# Bien dich va chay
!nvcc -o program program.cu
!./program
\end{lstlisting}

\section{Tài liệu tham khảo bổ sung}

\begin{itemize}
    \item NVIDIA CUDA C Programming Guide: \url{https://docs.nvidia.com/cuda/}
    \item CUDA By Example (Book): Sanders \& Kandrot
    \item Professional CUDA C Programming (Book): Cheng et al.
    \item Microsoft C\# Parallel Programming: \url{https://docs.microsoft.com/en-us/dotnet/standard/parallel-programming/}
    \item Thrust Library: \url{https://thrust.github.io/}
\end{itemize}
